\section{Starting Point of a Loop}
\label{sec:ll-loop-start}

\problemstatement{If a linked list has a cycle, return the node where the
cycle begins; otherwise return null. Use $O(1)$ extra space.}

\begin{patternbox}
After Floyd's slow/fast pointers meet inside the loop, a beautiful fact
kicks in: \textbf{resetting one pointer to the head and advancing both one
step at a time, they meet exactly at the loop's start}. It follows from the
distance arithmetic of the tortoise--hare.
\end{patternbox}

\subsection*{The distance argument}

Let the distance from head to the loop start be $L$, and let the slow and
fast pointers meet $k$ nodes into the loop. When they meet, the fast
pointer has travelled twice the slow pointer's distance; working through the
algebra shows that $L$ equals the remaining distance from the meeting point
around to the loop start. So if you place one pointer at the head and leave
the other at the meeting point, and move both one step at a time, they
converge precisely at the loop's entrance.

\begin{ideabox}[Head to Start Equals Meeting Point to Start]
The tortoise--hare meeting point is positioned so that the number of steps
from the head to the loop start equals the number of steps from the meeting
point (going forward around the loop) to the loop start. Two synchronized
one-step walks therefore collide at the start.
\end{ideabox}

\subsection*{Algorithm}

\begin{enumerate}
  \item Run slow/fast until they meet; if fast reaches null, return null (no
        loop).
  \item Move \texttt{slow} to the head. Advance \texttt{slow} and
        \texttt{fast} one step each until they are equal.
  \item That node is the loop's start.
\end{enumerate}

\subsection*{Dry run}

\dryrun{$1\to2\to3\to4\to2$ (loop starts at $2$).}

\begin{itemize}
  \item Slow/fast meet inside the loop; reset slow to head ($1$).
  \item One step each: they converge at node $2$ -- the loop start.
\end{itemize}

\subsection*{C++ implementation}

\begin{lstlisting}
#include <bits/stdc++.h>
using namespace std;

struct ListNode {
    int val; ListNode* next;
    ListNode(int v) : val(v), next(nullptr) {}
};

ListNode* detectCycle(ListNode* head) {
    ListNode* slow = head;
    ListNode* fast = head;
    while (fast && fast->next) {
        slow = slow->next;
        fast = fast->next->next;
        if (slow == fast) {                       // meeting point
            slow = head;
            while (slow != fast) { slow = slow->next; fast = fast->next; }
            return slow;                          // loop start
        }
    }
    return nullptr;
}

int main() {
    ListNode* a = new ListNode(1);
    a->next = new ListNode(2);
    a->next->next = new ListNode(3);
    a->next->next->next = new ListNode(4);
    a->next->next->next->next = a->next;          // 4 -> 2
    ListNode* start = detectCycle(a);
    cout << (start ? start->val : -1) << "\n";     // 2
    return 0;
}
\end{lstlisting}

\begin{complexitybox}
\textbf{Time Complexity.} $O(n)$. \textbf{Space Complexity.} $O(1)$.
\end{complexitybox}

\begin{takeawaybox}
The loop's start is found by resetting one pointer to the head after the
tortoise--hare collision and walking both one step at a time. The
underlying ``$L$ equals meeting-to-start'' identity is worth remembering --
it is asked as a follow-up to cycle detection constantly.
\end{takeawaybox}
